\section{Discussion and Limitations}
\label{sec:discussion}

\subsection{Six empirical boundaries}
\begin{table}[!t]
\centering
\small
\begin{tabular}{p{0.22\linewidth}p{0.48\linewidth}p{0.22\linewidth}}
\toprule
Boundary & Limit & Source \\
\midrule
Contamination ceiling
    & 20--25\% whale density; nothing robust works at $\geq$50\%
    & Appendix~\ref{app:whale} \\
Clean-data efficiency of robust nuisance
    & Huber($\delta=0.5$) costs 3$\times$ $\PEHE$ on IHDP
    & \S\ref{sec:experiments:ihdp} \\
Basis sensitivity
    & User supplies $\tx$ form; wrong basis $\to$ $\PEHE$ degradation
    & \S\ref{sec:experiments:tail_signal} \\
Compute
    & 5--20 s per fit (MCMC); not millisecond-latency territory
    & All benchmarks \\
Treatment type
    & Binary $W \in \{0, 1\}$; continuous dose unvalidated
    & Architectural \\
Sample size
    & Verified $200 \leq N \leq 5000$; very-large-$N$ extrapolated
    & Appendix~\ref{app:scaling} \\
\bottomrule
\end{tabular}
\caption{Six boundaries of the library's empirically-validated
envelope. Outside these regimes, use a different tool.}
\label{tab:boundaries}
\end{table}

\subsection{The limits of automated severity tuning}
Reviewers asked if we could provide an automated switch (e.g., via a robust
tail index) to toggle \texttt{severity} dynamically, avoiding the clean-data
efficiency cost when tails are light. We implemented and tested an automated
switch using the Hill estimator on preliminary residuals. While it successfully
assigned \texttt{severity=none} on clean data and \texttt{severe} under
5\% contamination, it failed catastrophically at 20\% contamination: dense,
shifted outliers trick the continuous Pareto-tail estimator into estimating
a light tail ($\hat{\alpha} \approx 6.0$), causing the pipeline to deactivate
robustification and incur massive bias ($>+6.0$). Therefore, we recommend
practitioners rely on domain knowledge to set \texttt{contamination\_severity}
rather than attempting to fully automate it.

\subsection{Scale sensitivity and MAD standardisation}
\label{sec:discussion:scale}
CatBoost's Huber-loss $\delta$ parameter operates on the raw outcome
scale, so the \texttt{contamination\_severity} API's
$\delta$-prescription is calibrated for unit-scale residuals. We
probe this on a synthetic dollar-scale DGP (true $\tau = 1000$,
$Y$ on order $\pm 10^4$, 10\% whales adding $+25{,}000$,
\texttt{severity="severe"} so $\delta = 0.5$ in CatBoost). Without
pre-standardisation, the estimator returns
$\hat\tau \approx 1085$ (8.5\% positive bias). With
$Y \mapsto Y/\widehat{\mathrm{MAD}}(Y)$ pre-standardisation
(then rescaling the posterior afterward), $\hat\tau \approx 943$
(5.7\% negative bias). Both are within $\pm 10\%$ of the truth, and
neither is catastrophic on this DGP. The bias direction flips
with rescaling, however --- a sign that the severity preset is
mildly miscalibrated for this scale. We recommend MAD
pre-standardisation as a default safety step on extreme-scale data;
where the bias of either side matters operationally, the
data-driven $(\delta, c)$ selector of
\S\ref{sec:experiments:other_paradigms} can be used for finer
tuning.

\subsection{When to use this library --- and when not}
The library is the right choice when the target is a calibrated
Bayesian posterior over $\tx$ under binary treatment, outcome data
is genuinely heavy-tailed or suspected to be, and second-level
compute (5--20 seconds per fit) is acceptable. It is particularly
well-suited to applications where the posterior itself is the object
of interest --- decision integration, subgroup contrasts, inequality
probabilities --- and not merely a point estimate with an error bar.
For continuous treatments, pure-prediction uplift tasks with no
uncertainty requirement, sub-second latency budgets, or very small
samples ($N < 200$), users should prefer causal forests, R-learners,
or TMLE variants as appropriate. We document these alternatives
explicitly rather than implicitly; the point is not that the
Bayesian X-Learner dominates but that it occupies a well-defined
niche the other tools do not cover.

\subsection{Broader impact and misuse risks}
\label{sec:discussion:risks}
The framework fills a practical gap for calibrated UQ on $\tx$ in
heavy-tailed domains, enabling principled decisions (e.g., subgroup
deployment, clinical go/no-go) where bootstrap CIs may underperform.
The tails-as-signal vs contamination guidance is particularly
helpful for practitioners. However, two risks deserve explicit
acknowledgement:

\textbf{Over-downweighting policy-relevant subgroups.}
If contamination is misdiagnosed (i.e., genuine tail heterogeneity
is treated as contamination via an overly aggressive severity
preset), the Welsch likelihood and Huber nuisance will
systematically downweight observations from rare but
policy-relevant subgroups, producing a posterior that appears
calibrated but is actually smoothing over real signal. Practitioners
should inspect the diagnostic residual plots
$\psi_W(r_i)$ vs $X_i$ (available via \texttt{plot\_diagnostics()})
to check whether downweighted residuals cluster in covariate
regions that might represent genuine subgroup effects.

\textbf{Overconfidence from mis-calibrated $\eta$.}
If $\eta$ is mis-calibrated (e.g., the ridge projection of
$\widehat{I}$ changes the target in a direction the practitioner
does not recognise) or if nuisance uncertainty is not propagated
(modular Bayes turned off when it should be on), the posterior
can appear tighter than it should be. We recommend two guards:
(i)~always report the cross-fit dispersion ratio $\rho$
(Table~\ref{tab:dispersion}) and turn on modular pooling when
$\rho > 0.15$; (ii)~inspect tail-residual plots of
$\psi_W(r_i)$ overlaid on the empirical distribution to verify
that the redescent is not suppressing structured signal.

\subsection{Honest failures and reviewer caveats}
Four observations we expect reviewers to probe. First, the IHDP
comparison of Table~\ref{tab:ihdp} is based on 5 replications,
which is the field's standard benchmark cadence but does not
separate method means at $\alpha = 0.05$ under Welch's test
(mean differences are within the replication-to-replication
standard error). The rank ordering is suggestive of relative
performance, not a conclusive superiority claim; the narrative of
$\PEHE \approx 0.56$ being a floor claim for the Bayesian machinery
should be read in that light. Second, the
clean-data efficiency cost analysed in
\S\ref{sec:experiments:efficiency} is not closable by relaxing
$\delta$: \texttt{contamination\_severity="none"} is the only
configuration that avoids it on IHDP, and switching loss types
is the only way to close the remaining gap. We see this as a
transparent cost the enum makes easy to navigate, not a defect of
the Huber framework. Third, coverage at very low whale density under
\texttt{severity="severe"} is zero because bias ($\approx 0.13$) is
comparable to CI half-width ($\approx 0.11$). The credible
intervals miss the true effect by a small amount, not by a
systematic failure of the posterior. A \texttt{severity="mild"}
default would widen intervals and recover coverage at this density
at the cost of worsening them at 20\%; the trade-off is documented
and explicit. Fourth, the \texttt{normalize\_extremes} path remains
in the codebase despite being empirically discouraged; it is
quarantined in a separate documented module
(Appendix~\ref{app:evt_path}) and disabled by default, but users
who pass \texttt{tail\_threshold} and \texttt{tail\_alpha}
explicitly can still activate it. We prefer transparent
architectural residue over silent removal that would break
prior experiments.

\subsection{Limitations beyond the empirical envelope}
\label{sec:discussion:limitations}
Three further limitations deserve explicit acknowledgement.
\textbf{(a) IHDP rank ordering is unstable across replications.}
At 5 reps the Bayesian X-Learner leads on mean PEHE; at 10 reps no
method statistically dominates ($p > 0.58$ for all pairwise Welch
tests against RX-Learner). Replications 6--10 are markedly harder
than 1--5 in the CEVAE preprocessing, which inflates every method's
mean. Our IHDP claim is therefore best read as ``competitive on
average, statistically indistinguishable from strong baselines at
the available sample size'' rather than as superiority. Extending
to the canonical 100 replications would clarify but is gated by
BART/BCF compute (minutes per replication).
\textbf{(b) BCF implementation parity.} Our BCF in
Table~\ref{tab:ihdp} is implemented in \texttt{pymc\_bart} with the
Hahn 2020 structural decomposition but does not fully replicate the
original tau-prior; SBCF~\citep{caron2022shrinkage} and
\texttt{stochtree}'s warm-start variants are not run. A reference
R-implementation of BCF would likely outperform our pymc-based BCF.
\textbf{(c) Hillstrom lacks ground-truth $\tau$.} Hillstrom is real
heavy-tailed outcome data, but per-unit treatment effects are
unobserved; we can only validate the marginal ATE there
(triangulated by placebo/symmetry checks in
\S\ref{sec:experiments:hillstrom}).
\textbf{(d) Auto-severity breaks at high contamination.} The Hill-
estimator-based severity selector
(\S\ref{sec:experiments:auto_severity}) routes light/moderate
contamination correctly but fails at 20\% whale density because the
non-robust pre-fit it relies on is itself contaminated; a robust
pre-fit pipeline would close the gap and is flagged as future work.
\textbf{(e) Theoretical calibration: partial.}
Proposition~\ref{prop:welsch_calibration} establishes an asymptotic
credible-interval calibration result for our Welsch generalised
posterior under the trace-formula learning rate
\eqref{eq:eta_a}--\eqref{eq:eta_tr}. The result inherits the standard generalised
Bernstein--von Mises regularity conditions and assumes a correctly
specified basis, the cross-fitting product-rate condition of
\citet{kennedy2020optimal}, and the bounded-influence property of
$\psi_W$. What remains open is a finite-sample, non-asymptotic
contraction rate (e.g., a PAC-Bayes-style result that quantifies
how (A1)--(A4) violations degrade coverage), and a corresponding
result for the modular-Bayes posterior of \S\ref{sec:experiments:nuisance_bootstrap} that propagates nuisance
uncertainty. The estimand-focused tuning of~\citet{alexopoulos2025rbci} gives
one concrete avenue: choose $(\eta, c)$ to optimise a proper
interval score for $\tau$ rather than the loss-likelihood-bootstrap
variance match we use in \S\ref{sec:experiments:learning_rate}.
Moreover, the ridge stabilisation of $\widehat{I}$ in the
$\eta$-estimation procedure introduces a controlled bias whose
magnitude depends on the ridge level $\lambda$ and the dimension $p$.
We report sensitivity to $\lambda \in \{0, 10^{-3}, 10^{-2}\}$
in \S\ref{sec:experiments:coverage} and find that
coverage is stable across $\lambda$ at $p \le 6$ but sensitive
at $p \ge 21$, where constrained eigenvalue projections (spectral
matching) or estimand-specific $\eta^\star(a)$
(\eqref{eq:eta_a}) should replace the trace formula.
\textbf{(f) Single semi-synthetic benchmark.} IHDP is the only
semi-synthetic CATE benchmark we report at scale. The ACIC 2016 /
2018 series and the Twins dataset
provide additional families of simulation conditions
that would tighten generality claims; we did not run them because the
data engineering (77 simulation conditions in ACIC 2016, 100
simulations each; Twins requires careful pre-processing of
mortality outcomes) is substantial and would distract from the
heavy-tail focus.
\textbf{(g.-1) GRF baselines not run.} Generalised random forests
\citep{athey2019generalized} with robustified splitting are a
canonical CATE alternative; we did not run them because the
reference implementation is the R \texttt{grf} package and we
chose to keep the experimental pipeline single-language (Python).
A multi-language reproduction is a natural extension.
\textbf{(g.0) Coverage replication count.} Most coverage tables
report 3 seeds per cell. The headline coverage findings (e.g.\ 100\%
under modular Bayes, 0\% under contaminated-normal at 20\% density)
have correspondingly wide binomial confidence intervals; we report
results without Wilson-style coverage error bars and recommend a
follow-up study with $\geq 30$ seeds per cell to substantiate
calibration claims at decision-grade precision. \textbf{(g) Robust-BART implementation incomplete.} We attempted a
contaminated-normal-residual T-BART variant in
\texttt{pymc\_bart} as a Bayesian-tree alternative to Welsch.
The mixture likelihood combined with BART's tree-splitting did not
sample successfully in our implementation (every replication
returned NaN). The Student-$t$ T-BART variant
(\S\ref{sec:experiments:ihdp}, $0.683 \pm 0.23$) succeeds and
is reported in Table~\ref{tab:ihdp}; an M-estimation-in-splitting
BART would require modifying the BART splitting rule, which
\texttt{pymc\_bart} does not expose.

\subsection{Practitioner defaults}
\label{sec:discussion:defaults}
A consolidated default-configuration table for users:

\begin{table}[!t]
\centering
\small
\begin{tabular}{p{0.27\linewidth}p{0.22\linewidth}p{0.43\linewidth}}
\toprule
Knob & Default & When to change \\
\midrule
\texttt{contamination\_severity} & \texttt{"none"} (XGB-MSE) & \texttt{"severe"} if outcome tails are visible (Hill index $\le 2$) or known whales \\
Welsch $c$ & 1.34 & 0.5 if you also set severity=severe \\
$\eta$-calibration mode & trace-formula \eqref{eq:eta_tr} & functional-specific \eqref{eq:eta_a} for a single contrast; RBCI for hard-coverage \\
$K$ (cross-fit folds) & 2 & no benefit from larger $K$ in our experiments \\
\texttt{use\_overlap} & False & True if min $\hat\pi \le 0.05$ or max $\hat\pi \ge 0.95$ \\
Modular Bayes & off (single cross-fit) & on, with $M = 8$, when dispersion ratio $\rho > 0.15$ \\
NUTS warmup / samples / chains & 400 / 800 / 2 & longer if $p \ge 50$ \\
$\beta$-prior scale $\sigma_\beta$ & 10.0 & 2.0 if $p \ge 10$ to avoid flat-region inefficiency \\
Student-$t$ prior $\nu$ & 3 & --- \\
\texttt{normalize\_extremes} & off (warning if activated) & avoid; use a tail-aware $\phi(x)$ basis instead \\
\bottomrule
\end{tabular}
\caption{Default configuration and adjustment guidance. Contamination
severity is the single most consequential knob; the dispersion-ratio
diagnostic of \S\ref{sec:experiments:coverage} tells the user
when modular pooling is needed.}
\label{tab:defaults}
\end{table}

\subsection{Configuration decision tree}
\label{sec:discussion:flowchart}
A linearised decision walk-through for users (referencing the
defaults table of \S\ref{sec:discussion:defaults}):

\begin{enumerate}[1.,leftmargin=*,itemsep=2pt]
\item \textbf{Inspect the outcome distribution.} If the empirical
Hill index $\hat\alpha \le 2$ on the upper tail of $|Y|$, contamination
is plausible. If $|Y|$ values span multiple orders of magnitude
(monetary outcomes, costs), enable
\texttt{normalize\_y\_for\_nuisance=True} (MAD-rescale).
\item \textbf{Inspect overlap.} If $\min \hat\pi(X) < 0.05$ or
$\max \hat\pi(X) > 0.95$, set \texttt{use\_overlap=True}.
\item \textbf{Set severity.} Default to \texttt{"none"}; switch to
\texttt{"severe"} if Hill index suggests heavy tails or domain
knowledge indicates whales / rare large events.
\item \textbf{Choose the basis $\phi(x)$.} For ATE only, use the
intercept. For tail-as-signal, use a tail indicator
$[1, \mathbf{1}(|x_j|>c)]$ at the suspected threshold $c$, or a
spike-and-slab over candidate thresholds when $c$ is unknown
(\S\ref{sec:experiments:basis_ablation}).
\item \textbf{Run the posterior.} If the dispersion ratio
$\rho > 0.15$ or you require nominal coverage, enable
\texttt{modular\_bayes=True} with $M = 8$. Otherwise default
single-cross-fit suffices.
\item \textbf{Calibrate $\eta$ for hard-coverage requirements.} The
trace-formula default is sharp; for guaranteed nominal coverage on
a target functional $a^\top \beta$, use $\eta^\star(a)$
(Proposition~\ref{prop:welsch_calibration}(a)) or RBCI
$\omega$-tuning.
\item \textbf{Diagnostic check.} If the smallest eigenvalue of
$\widehat I$ is near zero or the residual mass outside $|r| <
c/\sqrt{2}$ exceeds $0.6$, increase $c$ or enable overlap weights;
the library emits a warning in these regimes.
\end{enumerate}

\subsection{Future work}
\label{sec:discussion:future}
Five directions deserve attention. First, a proper Bayesian EVT
likelihood --- a generalised-Pareto tail mixture component
integrated with the Welsch bulk --- would promote the
tails-as-signal use case from ``supported via basis'' to
``supported via likelihood,'' with Hill estimation used for prior
elicitation rather than data rescaling. Preliminary results
(\S\ref{sec:experiments:hetero_evt_phase3}) show the concept
is viable but requires careful identifiability analysis. Second,
continuous treatments require a re-derivation of the Phase 2 DR
target and are scoped out of the current library; integration with
continuous-treatment DR-learners~\citep{kennedy2020optimal} is a
natural extension. Third, adaptive basis construction (e.g., via a
Gaussian-process prior over $\tx$) would relax the user's
responsibility to specify the functional form, at a compute cost
that may or may not be acceptable depending on application.
Fourth, a data-driven $\delta$ (and $c$) selector --- e.g., via
predictive interval scores or robust scale diagnostics --- would
replace the current fixed presets with an adaptive mechanism.
The tail-index-based auto-severity selector of
\S\ref{sec:experiments:auto_severity} is a first step, but
its failure at high contamination (\S\ref{sec:experiments:auto_severity})
motivates a robust-pre-fit variant that uses a preliminary
Huber-fitted S-learner before tail estimation.
Fifth, a proper generalised-Bayes
calibration analysis of the Welsch pseudo-likelihood
\citep{bissiri2016general,lyddon2019general}, possibly with a
decision-theoretic $\omega$-selector in the style of
\citet{alexopoulos2025rbci}, would replace the empirical
calibration argument of \S\ref{sec:method:phase3} with a
formal one and enable joint optimisation of $(\eta, c)$ for
the causal estimand of interest.
